\chapter{Тестирование и анализ результатов}
\label{ch:testing}

\section{Цели и методика тестирования}

Тестирование направлено на проверку:
\begin{enumerate}
  \item корректности REST API и валидации входных данных;
  \item работоспособности алгоритма планирования на типовых сценариях;
  \item устойчивости к недоступности OSRM (fallback);
  \item соответствия метрик ожидаемому поведению при изменении предпочтений.
\end{enumerate}

Использованы методы: ручное функциональное тестирование через UI и API,
анализ ответов JSON, сравнение маршрутов при различных векторах $\mathbf{w}$.

\section{Тестовое окружение}

\begin{table}[H]
\centering
\caption{Конфигурация тестового окружения}
\label{tab:test-env}
\begin{tabular}{@{}ll@{}}
\toprule
\textbf{Параметр} & \textbf{Значение} \\
\midrule
ОС & Windows 10 / Ubuntu 22.04 \\
Node.js & 18 LTS \\
Backend URL & http://127.0.0.1:3000 \\
Frontend URL & http://127.0.0.1:5173 \\
OSRM & router.project-osrm.org \\
Демо-регион & Центр Москвы, 15 POI \\
\bottomrule
\end{tabular}
\end{table}

Запуск: \texttt{npm run dev} из корня проекта или \texttt{scripts/dev.ps1}.

\section{Функциональные тест-кейсы}

\begin{table}[H]
\centering
\caption{Результаты функционального тестирования}
\label{tab:test-cases}
\small
\begin{tabular}{@{}p{1.2cm}p{4.2cm}p{4.5cm}p{2.5cm}@{}}
\toprule
\textbf{ID} & \textbf{Действие} & \textbf{Ожидаемый результат} & \textbf{Статус} \\
\midrule
TC-01 & GET /api/health & \texttt{\{status: 'ok'\}} & Пройден \\
TC-02 & GET /api/poi & 15 объектов & Пройден \\
TC-03 & POST plan без body & 400 Bad Request & Пройден \\
TC-04 & POST plan, lat=999 & 400 (валидация) & Пройден \\
TC-05 & Plan: дефолтные prefs & Маршрут, score $>$ 0 & Пройден \\
TC-06 & parks=10, остальное=0 & Маршрут через парки & Пройден \\
TC-07 & waterfront=10 & Набережные в highlights & Пройден \\
TC-08 & Клик карты: старт/финиш & Координаты обновлены & Пройден \\
TC-09 & Сохранённые маршруты & Список после plan & Пройден \\
TC-10 & DELETE /api/routes/:id & 404 после удаления & Пройден \\
\bottomrule
\end{tabular}
\end{table}

\section{Анализ алгоритма на примере}

\textbf{Сценарий:} старт --- Красная площадь $(55.7558, 37.6173)$, финиш --- Парк
Горького $(55.7293, 37.6017)$, предпочтения: parks=9, greenZones=7, airQuality=8,
остальные --- средние.

\textbf{Ожидаемое поведение:}
\begin{itemize}
  \item высокий $\mathrm{score}$ у Парка Горького, Музеона, Нескучного сада;
  \item Dijkstra включает 2--4 POI на пути;
  \item OSRM строит polyline по аллеям, не через Москва-реку «напрямую»;
  \item $\mathrm{detourRatio} > 1$, но ограничен разумными рамками;
  \item $\mathrm{Score}$ в ответе API --- 55--85 (зависит от набора POI).
\end{itemize}

\textbf{Наблюдаемые \texttt{graphStats}:} типично nodes = 4--10, edges = 6--20,
pathWeight --- сумма весов \eqref{eq:edge-weight} по найденному пути.

\section{Влияние вектора предпочтений}

Проведено сравнение трёх профилей при фиксированных $S$ и $E$:

\begin{table}[H]
\centering
\caption{Сравнение маршрутов при различных предпочтениях}
\label{tab:pref-comparison}
\begin{tabular}{@{}lccc@{}}
\toprule
\textbf{Метрика} & \textbf{«Парки»} & \textbf{«Набережные»} & \textbf{«Велодорожки»} \\
\midrule
parks / waterfront / bike & 10 / 0 / 0 & 0 / 10 / 0 & 0 / 0 / 10 \\
Типичные highlights & Горького, Музеон & Крымская наб. & Велодорожки \\
greenCoverage, \% & 60--85 & 40--60 & 20--40 \\
Score & 65--80 & 55--70 & 45--60 \\
\bottomrule
\end{tabular}
\end{table}

Вывод: модель \eqref{eq:poi-score}--\eqref{eq:edge-weight} обеспечивает
различимость маршрутов --- система реагирует на изменение $\mathbf{w}$ предсказуемым
образом.

\section{Тест fallback OSRM}

При блокировке OSRM (неверный URL в \texttt{.env}) система:
\begin{itemize}
  \item возвращает \texttt{routingSource: 'direct'};
  \item waypoints --- вершины графа Dijkstra;
  \item score снижается до $\approx$30--35;
  \item UI отображает предупреждение о прямой линии.
\end{itemize}

Graceful degradation подтверждён --- приложение не падает с необработанным
исключением.

\section{Производительность}

При $|V| \leq 10$ этап Dijkstra занимает $< 1$~мс. Доминирует latency OSRM
(1--8~с в зависимости от числа waypoints и нагрузки публичного сервера). Общее
время POST \texttt{/api/routes/plan} --- 2--15~с (ограничено timeout 15~с).

\section{Выявленные ограничения}

\begin{enumerate}
  \item Метрики greenCoverage и bikePath --- эвристики, не GIS-анализ;
  \item При равных score порядок POI может зависеть от порядка в массиве;
  \item Публичный OSRM не гарантирует SLA для production;
  \item In-memory storage не подходит для многопользовательского режима.
\end{enumerate}

\section{Выводы по главе}

Функциональное тестирование подтвердило работоспособность всех заявленных
возможностей. Алгоритм демонстрирует корректную реакцию на пользовательские
предпочтения. Выявленные ограничения не препятствуют использованию системы как
учебного прототипа и базы для дальнейшего развития.
